\section{Method}\label{sec:method}
\paragraph{Speaker prompt.} A's system prompt contains the memories, the secret, the cover topic and the list of banned names. It also lists which memories mention the helper and which mention the venue
(a deterministic function of information A already holds) and gives one example of the hinting style. The prompt is identical for every speaker condition.
\paragraph{Listeners.} B and C are the same instruction-tuned model (\x{bigGen}), used zero-shot. B reads the memories and the post-switch chat. C reads only the post-switch chat, the rules of the game and three solved training examples.
Each listener yields a distribution over the 10 helpers and 10 venues by scoring each option as the continuation of ``Helper:'' or ``Venue:'' (log-probability of the first branching token, normalised over options).
We first trained 1.5B LoRA listeners on the warm-start transcripts; they did not exceed a guessing baseline (Appendix~\ref{app:trained}), so all reported rewards and evaluations use the prompted pair.
\paragraph{Warm start.} The 14B model plays A (with the secret) and B (without it); A turns that name any inventory item are re-sampled up to eight times (warm start only).
A conversation is \emph{good} if A does not leak, the sense judge passes, the prompted B recovers both slots and the prompted C at most one. Supervised fine-tuning on the \x{sftExamples} A turns of the \x{warmGood} good conversations gives the \emph{sft} adapter (condition S1), which is also the KL reference.
\paragraph{Reward.} For one rollout, $r=-1$ if any A turn names a person or place, or if the sense judge's score is below $\tau$; otherwise $r=\tfrac12(p_B^{h}+p_B^{v})-\lambda\,\tfrac12(p_C^{h}+p_C^{v})$, where $p$ is the probability a listener assigns to the true value and $\lambda=1$.
The sense judge is the base 3B model with adapters disabled, scoring P(Yes) to ``Is this a natural, coherent conversation\ldots?''; $\tau=\x{tau}$ was chosen on \x{judgeNGood} good and \x{judgeNBad} corrupted warm-start conversations (turns swapped from another episode, or one turn replaced by random words); at this threshold it kept \x{judgeKeepGood}\% of the good ones and rejected \x{judgeRejectBad}\% of the corrupted ones. The good conversations had already passed the judge at the default threshold of 0.5, so the keep rate is optimistic.
\paragraph{Multi-turn GRPO.} For $P$ episodes we sample $G$ rollouts each: A turns come from the policy adapter, B turns from the frozen base model, for four turns (A, B, A, B). B turns that name an inventory item are re-sampled up to three times; if the last sample still names one it is kept, and the reward does not penalise it (Section~\ref{sec:analysis}).
Advantages are group-normalised, $A_i=(r_i-\bar r)/(\sigma+10^{-4})$, and zero-variance groups contribute nothing. The loss on A-turn tokens only is a PPO-style clipped surrogate ($\epsilon=0.2$) plus $\beta\,\mathrm{KL}$ to the sft reference (k3 estimator), averaged over tokens within a sequence and then over sequences; both A turns of a rollout share its advantage.
\paragraph{Evaluation.} For each speaker condition we generate one conversation per test episode with identical decoding settings, B speaker, judge and leak checker, and score it with the prompted B and C. As a check that the speaker did not exploit the particular outsider prompt it was trained against, we also score every test conversation with a held-out outsider C$_\text{alt}$: the same model with reworded instructions and three solved examples from different training episodes, never used in the reward.
